\section{Conclusion}
\label{sec:conclusion}

We summarise the state of Pollard rho on current NVIDIA hardware as
measured in this work.

\paragraph{Per-card throughput.} In 2010 the
best GTX~295 kernel ran ECC2K-130 at $0.064$ billion iterations per
second. On one RTX~PRO~6000 Blackwell the same walk runs at $14.64$ in
the audited campaign build and $15.88$ in the best later build, and the
whole $2^{60.9}$ budget is about five GPU-years. Roughly a sixth of that
factor is a bigger, faster chip; the rest is arithmetic. Sixteen paired
steps took the client from a $0.85$ bitsliced control to $14.64$, and
three representation changes did most of the work: Frobenius
permutation networks, the native carry-less product, and the smaller
batch that the native product made possible.

\paragraph{Time to solution.} At $15.88$ billion
iterations per second, five RTX~PRO~6000s expect to find the collision
in about a year, 56 in a month, 290 in a week, and a thousand in about
three days, half of them spent waiting for the last trail to reach a
distinguished point at cutoff 32. The work costs about \$114k at list
price and \$38k--\$56k on the spot quotes we saw, whatever the fleet
size. The same rates leave P-256 about $6\cdot10^{21}$ card-years away
(Section~\ref{sec:solvetime}).

\paragraph{Carry-less multiplication across dies.}
PTX~9.3 gives every part from \texttt{sm\_80} onward a carry-less
multiply, but not at the same rate. On the RTX~PRO~6000 it issues at
about $1.7$ lanes per SM-clock, the cost of $38$ logic instructions; on
the B200 and H100 it is ten to eighteen times faster. The workstation
part still wins on this walk, per card, per SM and per dollar, because
the walk is bound by integer logic and the RTX~PRO~6000 has the most
SMs at the highest clock for the money. On the data-centre parts the
logic pipe binds instead, the right kernel moves work onto the
carry-less unit, and kernel tuning does not travel between them. For
any walk with one affine addition per step, $28$ billion iterations per
second on one RTX~PRO~6000 is out of reach on both pipes.

\paragraph{Choice of iteration function.} The table walk
does the arithmetic of a step more cheaply, and reached $20.08$ under
cycle rule v2. But v2's exits from fruitless cycles depended on how a
trail arrived, and the exact v3 rule that fixes this runs at $5.10$
today. Closing that gap, with a deferred or off-critical-path cycle
check, is the largest open optimisation in this paper, and even then a
table walk must win on cost per solve, which nobody has measured.

\paragraph{Transfer to prime fields.} On NIST P-256 and P-384 on the same card, explicit carry
chains, a wider shared inversion and the same refusal of register-cap
trades gave $1.52\times$ and $2.61\times$, to $1.632$ and $0.548$ billion
iterations per second. The batch width, a modest knob on the binary
walk, was the biggest lever on the prime one, because a Fermat
inversion costs some forty times as many multiplications as a
normal-basis one. A P-256 step remains about nine times as expensive
as an ECC2K-130 step on the same card, which is the carry-less unit's
advantage in one number.

\paragraph{Remaining headroom.} None of the speedup
came from changing the algorithm: every step kept Bailey et al.'s walk
and its distinguished points bit for bit, and the gains split into
representation and arithmetic ($6.3\times$), batch and geometry
($1.5\times$), memory layout ($1.3\times$) and the new instruction
($1.2\times$). The same pipe measurements bound what is left. Any walk with
one affine addition per step tops out at about $24$--$26$ billion
iterations per second on one RTX~PRO~6000, so the best $\sigma$-walk has
roughly $1.5\times$ to go, all of it in integer work; a B200 would need a
$1.55\times$ cut in logic instructions to reach $30$; the F2 FPGA stops near
$9.3$. On CPUs the instruction sets matter as much as on GPUs: three-input
logic and wider words give $2.4$--$3\times$ to the bitsliced client, and a
packed client on the CPU's own carry-less multiply runs $2.9\times$ faster
than with a software product.

\paragraph{Asymptotic cost and the campaign.} The campaign's walk needs $2^{60.91}$--$2^{60.92}$ iterations,
which agrees with the 2009 estimate. Index calculus on the curve costs
$10^7$ times as much as rho, and the best measured decomposition is
$35$--$51$ bits short of its budget. The campaign stored $486$ million
points and $2^{56.68}$ iterations, $5.4\%$ of the expected work, with no
collision, every record re-walkable from its seed; its fleet was
reclaimed on 2026-09-18 and collection is paused. Whether the remaining
$94.6\%$ is walked is a question of money, not mathematics, and at five
GPU-years it is a much smaller question than it was in 2009.
